\cotefile{67}
\scanpage{53}
% slanted
\margin{torseurs sous groupe \unsure{diédrique|orthogonal|toriques} $\longleftrightarrow$ \deleted{(\unsure{symétrie, 1 tourne} $\unsure{\mathbb U/\mathbb Z_{2}}$, \unsure{à 2 élém.|const. échec})}}

\diagram{\begin{tikzcd}[ampersand replacement=\&,column sep=small,row sep=1em]
\text{Cercles toriques} \arrow[r,leftrightarrow] \&
\substack{\text{\deleted{\unsure{espace}} plan vectoriel \inserted{réel}}\\ \text{euclidien \deleted{\unsure{orié}}}} \arrow[r,leftrightarrow] \&
\left[\substack{\text{Droites vectorielles \inserted{unitaires} sur corps }C\\\text{iso. à }\mathbb C\ \text{(mais pas can. !)}}\right]\\
\substack{\text{Cercles toriques}\\\text{orientés}} \arrow[r,leftrightarrow] \arrow[d,leftrightarrow] \&
\substack{\text{\deleted{\unsure{espaces}} plans vectoriel}\\\text{réel orienté}} \arrow[r,leftrightarrow] \&
\substack{\text{\deleted{plans} \inserted{droites vectorielles} complexes }\deleted{\unsure{D}}\,T\\
\text{\deleted{structure \unsure{euclidienne}} \inserted{unitaires}}\\
\text{i.e. \deleted{avec} iso } T\simeq\overline{T^\vee}\\
T^\vee\simeq\overline T\ \text{i.e.}}\\
\text{torseurs sous }\mathbb U \arrow[rr,leftrightarrow] \& \&
\substack{\text{droites vectorielles complexes}\\ \text{conformes, avec \deleted{\unsure{disques}} jauge}\\
\text{\deleted{conformes} \deleted{\unsure{vectoriels}} conforme fixée }D}
\end{tikzcd}}

\deleted{\unsure{Le degré|Le topos|Le groupe} des objets des catégories}

\deleted{Soit $C', C''$ deux cercles toriques orientés, \unsure{alors} on a
$\operatorname{Isom}(C'^{-1},C'')\simeq\operatorname{Isom}(C''^{-1},C')$}

% slanted
\margin{anti-isom. des $\mathbb R^{*+}$-torseurs}

\uline{NB} Pour $D$ droite complexe fixée, on a Ens des jauges sur $D$
$=D^{\unsure{*|\times}}/\mathbb U$, et $J(D)\simeq J(D^\vee)_{\unsure{\varepsilon}}$ (car \(=\) groupe de rotations $\mathbb U_{\varepsilon}$).

Le passage au cercle torique \inserted{orienté} dual correspond au passage à l'espace vectoriel \deleted{\unsure{réel}} \deleted{\unsure{euclidien}} \inserted{complexe dual} $\deleted{\unsure{T}}T^\vee$ ($\simeq\overline T$)

La donnée, pour deux \deleted{\unsure{espaces}} \inserted{droites} vect complexes $T',T''$, d'un système $((D',D''),u)$, où $D',D''$ sont des jauges des $T',T''$ et $u\in\partial D'\mathbin{\wedge}_{\mathbb U}\partial D''$.
\[
\begin{aligned}
(&\simeq\operatorname{Isom}_{\mathbb U}\,\partial D'^{-1},\partial D''\\
 &\simeq\operatorname{Isom}_{\mathbb U}(\partial D''^{-1},\partial D')
 \simeq\operatorname{Isom}(T'^{-1},T'')\simeq\operatorname{Isom}(T''^{-1},T')\simeq\mathrm{Bases}(T'\otimes T'')\\
 &\simeq T'^{*}\mathbin{\wedge}_{\mathbb C^{*}}T''^{*}\}\simeq
 \deleted{\unsure{f\ ;D',D''\ ;\ \text{non}}}\ \text{dualités entre }(T',T'')\ )
\end{aligned}
\]
revient à la donnée d'un $v\in(T'\otimes_{C}T'')^{*}$ i.e. une dualité entre $T',T''$, \uline{plus} une jauge \inserted{$D'$} dans $T'$ (ou ce qui revient au $=$ \unsure{grâce à|grâce} \unsure{v}, plus une jauge $D''$ dans $T''$)
% archive sheet number: 55
